\chapter{Sampling textures}
\label{ch:sampling}

A shader asks for the value of a texture at an arbitrary point $(u,v)$, which almost never
falls on a texel centre. How the answer is computed---\emph{texture filtering}---affects both
the quality of the image and, in Texture Explorer, the agreement between the numbers the
probe computes on the CPU and the pixels the GPU draws.

\section{Addressing}
Coordinates outside $[0,1]$ are folded back according to an \emph{addressing mode}. Texture
Explorer uses \emph{wrap} (repeat): the texture tiles the plane, and texel indices are taken
modulo the size of the image. In C++, where the remainder of a negative number is negative,
the correct index is
\[
  \operatorname{wrap}(i,n) = ((i \bmod n) + n) \bmod n .
\]
Wrapping matters even inside $[0,1]$: the texels on the left border must be interpolated with
those on the right border, which is what makes tiling textures seamless.

\section{Nearest and bilinear filtering}

\begin{figure}[ht]
\centering
\begin{tikzpicture}[scale=2.2]
  \draw[black!25] (-0.5,0.5) grid[step=1] (2.5,-1.5);
  \foreach \x/\y/\n in {0/0/{h_{00}},1/0/{h_{10}},0/-1/{h_{01}},1/-1/{h_{11}}} {
    \fill (\x+0.5,\y-0.5) circle (1.6pt);
    \node[above left,font=\small] at (\x+0.5,\y-0.5) {$\n$};
  }
  \fill[npcolor] (0.85,-0.85) circle (1.8pt);
  \node[npcolor,below right,font=\small] at (0.85,-0.85) {$(x,y)$};
  \draw[|<->|,tcolor] (0.5,-0.3) -- node[above,font=\small] {$f_x$} (0.85,-0.3);
  \draw[|<->|,bcolor] (0.3,-0.5) -- node[left,font=\small] {$f_y$} (0.3,-0.85);
  \draw[dashed,black!50] (0.85,-0.5) -- (0.85,-1.5);
  \draw[dashed,black!50] (0.5,-0.85) -- (1.5,-0.85);
\end{tikzpicture}
\caption{Bilinear filtering interpolates the four texel centres that surround the sample
point, with weights given by the fractional offsets $f_x$, $f_y$.}
\label{fig:bilinear}
\end{figure}

\emph{Nearest} filtering returns the texel that contains the point; magnified, the texture
looks blocky. \emph{Bilinear} filtering interpolates between the four nearest texel centres.
In texel units the point is at $x = uW - \frac12$, $y = vH - \frac12$ (the $-\frac12$ moves
the origin to the first texel's centre); with $x_0=\lfloor x\rfloor$, $y_0=\lfloor y\rfloor$
and the fractions $f_x = x - x_0$, $f_y = y - y_0$,
\[
  h(x,y) = (1-f_y)\bigl[(1-f_x)\,h_{00} + f_x\,h_{10}\bigr] + f_y\bigl[(1-f_x)\,h_{01} + f_x\,h_{11}\bigr],
\]
where $h_{ab}$ is the texel $(x_0+a,\ y_0+b)$ (Figure~\ref{fig:bilinear}). Texture Explorer
implements exactly this formula on the CPU (Section~\ref{sec:cpusampling}), so that the probe
reads the same values as the GPU. The result is continuous but its derivative is not: it
jumps at texel centres, which will matter when we differentiate height maps.

\section{Minification and mipmaps}
When the texture is seen from far away, one pixel of the screen covers many texels. Sampling
just one point of that footprint is \emph{aliasing}: fine detail turns into noise that
shimmers when the camera moves. The standard remedy, \emph{mipmapping} \cite{williams1983},
precomputes a pyramid of images, each half the size of the previous one and obtained by
averaging $2\times2$ blocks, and chooses for each pixel the level whose texels are about the
size of the pixel. If $\rho$ is the number of texels of level~0 that the pixel spans, measured
with the screen-space derivatives of $(uW, vH)$,
\[
  \rho = \max\Bigl(\Bigl\lVert\dd{(uW,vH)}{x}\Bigr\rVert,\ \Bigl\lVert\dd{(uW,vH)}{y}\Bigr\rVert\Bigr),
  \qquad \lambda = \log_2 \rho ,
\]
the GPU samples level $\lambda$, interpolating between the two nearest levels
(\emph{trilinear} filtering). When the footprint is long and thin---a surface seen at a grazing
angle---a single level is either too blurry along one direction or aliased along the other;
\emph{anisotropic} filtering takes several samples along the long axis. Texture Explorer
creates every texture with its complete mipmap chain (generated by the GPU) and uses a
sampler with $8\times$ anisotropic filtering.

\paragraph{Mipmaps in a vertex shader.} The level $\lambda$ is computed from the difference
between neighbouring \emph{pixels}, which only exist in the pixel shader. A vertex shader that
reads a texture---as displacement mapping does---must choose the level itself. The right
choice depends on how densely the mesh samples the texture: Texture Explorer's meshes have
$256\times256$ cells per patch and its maps $1024\times1024$ texels, so each cell spans four
texels and the vertex shader reads level $\log_2(1024/256)=2$. Reading level~0 instead would
make the displaced surface jitter as tiny height details fall between vertices.

\paragraph{Derivatives of a texture.} The bump-mapping code of Chapter~\ref{ch:height}
samples the bump map one texel to each side of the point, always at level~0, so that the CPU
and the GPU compute the same slope. This is simple and exact but aliases when the solid is
small on the screen; a production renderer would scale the offsets with the mip level, or use
a normal map with its own mipmaps.

\section{Colour spaces}
\label{sec:srgb}

The values stored in an 8-bit colour image are not proportional to light. They are encoded with
the \emph{sRGB transfer function}, which spends more of the 256 codes on dark tones, where the
eye is more sensitive. Decoding a stored value $C\in[0,1]$ to linear light gives
\[
  C_{\mathrm{lin}} =
  \begin{cases}
    C/12.92, & C\le 0.04045,\\[2pt]
    \bigl((C+0.055)/1.055\bigr)^{2.4}, & C>0.04045,
  \end{cases}
\]
which is close to the power law $C^{2.2}$ (Figure~\ref{fig:srgb}).

\begin{figure}[ht]
\centering
\begin{tikzpicture}
\begin{axis}[width=7.5cm,height=5.5cm,xmin=0,xmax=1,ymin=0,ymax=1,
  xlabel={stored value $C$},ylabel={linear light},legend pos=north west,
  legend style={font=\small},grid=major,grid style={black!10}]
  \addplot[domain=0:1,samples=60,thick,npcolor] {x<=0.04045 ? x/12.92 : ((x+0.055)/1.055)^2.4};
  \addlegendentry{sRGB}
  \addplot[domain=0:1,samples=60,dashed,tcolor] {x^2.2};
  \addlegendentry{$C^{2.2}$}
  \addplot[domain=0:1,samples=2,dotted,black] {x};
  \addlegendentry{identity}
\end{axis}
\end{tikzpicture}
\caption{The sRGB decoding curve and its usual approximation.}
\label{fig:srgb}
\end{figure}

Lighting computations add and multiply amounts of light, so they must work with linear values.
Hence the rules that Texture Explorer follows:
\begin{itemize}
\item The \emph{colour} map is stored in sRGB and must be decoded before lighting. The GPU does
it for free when the texture is read through a view with an \texttt{\_SRGB} format.
\item \emph{Data} maps---bump, displacement, roughness, normal---store numbers, not colours,
and must \emph{not} be decoded. They are read through a plain \texttt{UNORM} view, which
returns the stored byte divided by $255$.
\item The scene is rendered into an sRGB target: the GPU encodes the linear result when it
writes each pixel, and also when it averages the samples of multisample antialiasing, which is
therefore done correctly in linear light.
\item The user interface, which works with display values, reads both the textures and the
rendered scene through plain views, so that the stored values reach the screen untouched.
\end{itemize}
One texture, two interpretations: Direct3D allows this when the texture is created with a
\emph{typeless} format (\texttt{R8G8B8A8\_TYPELESS}) and each \emph{view} specifies how to read
it (Section~\ref{sec:typeless}).
